\textbf{结论名称：} 中线长公式（阿波罗尼斯定理）

\textbf{适用条件：} 任意三角形 $ABC$，已知三边长 $a,b,c$，$a$ 边所对的中线 $m_a$。

\textbf{变量说明：} $a,b,c$ 分别为 $\triangle ABC$ 的顶点 $A,B,C$ 所对边的边长，$m_a$ 为边 $a$ 上的中线长，其余中线 $m_b,m_c$ 类推。

\textbf{核心公式：}
\[
\boxed{m_a^{2}=\frac{2b^{2}+2c^{2}-a^{2}}{4}}
\]
对称地，$m_b^{2}=\dfrac{2a^{2}+2c^{2}-b^{2}}{4}$，$m_c^{2}=\dfrac{2a^{2}+2b^{2}-c^{2}}{4}$。

\textbf{使用场景：} 已知三角形三边求中线长；判断中线与边的关系；在解三角形、向量或几何证明中快速计算中线平方。

\textbf{简单例子：} 设 $\triangle ABC$ 中 $a=6$，$b=5$，$c=5$，求边 $a$ 上的中线长 $m_a$。
解：直接代入公式
\[
m_a^{2}=\frac{2\times 5^{2}+2\times 5^{2}-6^{2}}{4}= \frac{50+50-36}{4}=\frac{64}{4}=16,
\]
故 $m_a=4$。